\section{Method}
\label{sec:method}

\subsection{A residency model}

Let \(P\) be the set of hash partitions and \(s_p\) the resident size of
partition \(p \in P\). A probe batch \(b\) targets exactly one partition,
written \(\pi(b)\). For a dispatch order \(\sigma\) over a buffer of \(n\)
batches, the reuse distance of the \(i\)-th batch is the total resident
size touched since the previous batch with the same target:

\begin{equation}
\label{eq:reuse}
d_\sigma(i) = \sum_{j \in R_\sigma(i)} s_{\pi(\sigma_j)},
\qquad
R_\sigma(i) = \{\, j < i : \pi(\sigma_j) \neq \pi(\sigma_i) \,\}.
\end{equation}

A batch hits cache when \(d_\sigma(i) < C\), the effective cache capacity.
Minimizing misses is therefore minimizing the number of indices whose
reuse distance exceeds \(C\), which grouping by target achieves exactly:
consecutive batches with equal targets have reuse distance zero.

\subsection{The scheduler}

The online scheduler holds a bounded buffer of \(n\) batches and always
dispatches the group with the largest count, breaking ties by arrival
time so that no batch can be starved. The bound on \(n\) is what keeps
latency predictable:

\begin{itemize}
\item A batch waits at most \(n - 1\) dispatch slots.
\item Output order is restored by the operator's existing sequence
  numbers, so no downstream operator observes the reordering.
\end{itemize}

\begin{table}[htbp]
\centering
\begin{tabular}{lrr}
\toprule
Buffer \(n\) & Miss rate & Throughput \\
\midrule
1 (arrival order) & 0.41 & 1.00 \\
8 & 0.22 & 1.44 \\
32 & 0.13 & 1.71 \\
128 & 0.11 & 1.74 \\
\bottomrule
\end{tabular}
\caption{Buffer size against measured miss rate and normalized
throughput on the \textsc{tpch-q3} workload.}
\label{tab:buffer}
\end{table}

Table~\ref{tab:buffer} shows that the returns flatten by \(n = 32\),
which is the value used throughout Section~\ref{sec:results}.
